\documentclass[11pt,a4paper]{article}

\usepackage[T1]{fontenc}
\usepackage[utf8]{inputenc}
\usepackage[danish]{babel}
\usepackage{mathpazo}
\usepackage[scaled=0.88]{beramono}
\usepackage[a4paper,margin=3.1cm]{geometry}
\usepackage{booktabs,tabularx}
\usepackage{graphicx}
\usepackage{amsmath,amssymb}
\usepackage[hang,flushmargin,bottom]{footmisc}
\usepackage[font=small,labelfont=sc,labelsep=period]{caption}
\usepackage{setspace}
\usepackage{fancyhdr}
\usepackage{titlesec}
\usepackage[hidelinks]{hyperref}

\titleformat{\section}{\normalsize\scshape\bfseries}{\thesection}{0.8em}{}
\titleformat{\subsection}{\normalsize\itshape}{\thesubsection}{0.8em}{}
\titlespacing*{\section}{0pt}{1.6em}{0.6em}
\setlength{\parindent}{1.5em}
\setlength{\parskip}{0pt}
\onehalfspacing

\pagestyle{fancy}\fancyhf{}
\renewcommand{\headrulewidth}{0pt}
\fancyfoot[C]{\small\thepage}

% Marker alt der mangler et tal, saa det kan findes med grep -rn "TBD"
\newcommand{\tbd}[1]{\textbf{[TBD: #1]}}

\fancyhead[R]{\footnotesize\scshape Timing luck}

\begin{document}
\selectlanguage{danish}
\thispagestyle{empty}

\begin{center}
  {\Large\bfseries Hvor meget af en regimestrategis afkast er held med datoen?\par}
  \vspace{0.3em}
  {\large Timing luck i månedligt rebalancerede aktieallokeringer\par}
  \vspace{1.4em}
  {\large Albert \textsc{[Efternavn]}\footnote{Økonomisk Institut,
   Københavns Universitet. Kontakt: \texttt{[mail]}. Replikationskode og data:
   \url{https://github.com/[bruger]/08-timing-luck}.}\par}
  \vspace{0.8em}
  {\normalsize Denne version: \today\par}
\end{center}

\vspace{1.2em}
\begin{center}\begin{minipage}{0.85\textwidth}\small
\begin{center}\textsc{Resumé}\end{center}
\noindent
En strategi, der skifter mellem offensiv og defensiv aktieeksponering én gang
om måneden, er i virkeligheden 21 forskellige strategier --- én for hver
mulige handelsdag. Jeg kører den samme regimemodel med alle 21
rebalanceringsdage på amerikanske data fra 2003 til 2026. Sharpe varierer fra
0,618 til 0,735 alene som følge af datovalget, et spænd på 0,117. Til
sammenligning slår medianvarianten en passiv aktieposition med 0,096 i Sharpe.
Spredningen på tværs af arbitrære datovalg overstiger altså den effekt,
strategien skal påvise. Opdelt på episoder er billedet skarpere: i finanskrisen
går det maksimale tab fra $-9{,}3$ til $-24{,}3$ procent alt efter
rebalanceringsdag, mens modellen i marts 2020 gav nul beskyttelse og i 2022 var
marginalt dårligere end en passiv position. Hele risikoreduktionen stammer fra
én episode, og dens størrelse er i vid udstrækning tilfældig.
\end{minipage}\end{center}

\vspace{0.8em}
\begin{center}\begin{minipage}{0.85\textwidth}\small
\noindent\textsc{Nøgleord:} regimeskift, rebalancering, timing luck, aktieallokering\\
\noindent\textsc{JEL:} G11, G17
\end{minipage}\end{center}

\vfill
\begin{center}\small\textit{Studieprojekt på offentlige data. Ikke investeringsrådgivning.}\end{center}
\clearpage\setcounter{page}{1}

\section{Motivation}
Backtests af månedlige strategier rapporteres næsten altid for én
rebalanceringsdato, typisk månedens sidste handelsdag. Valget er arbitrært, og
hvis resultatet afhænger stærkt af det, er den rapporterede præstation delvis
held. Det er en test, der koster lidt kode og kan ændre fortolkningen af en hel
strategi.

Jeg finder, at spredningen i Sharpe på tværs af de 21 mulige
rebalanceringsdage er større end strategiens fortrin over en passiv
benchmark. Den centrale observation er dog ikke spredningen i sig selv, men
hvor den kommer fra: modellen leverer sin risikoreduktion i præcis én af tre
store nedture, og netop i den episode er udfaldet mest følsomt over for
datovalget.

\section{Data}
Daglige afkast på SPY (bredt amerikansk aktieindeks) som offensiv position og
SHY (korte statsobligationer) som defensiv, udbyttejusterede, fra
2.~januar 2003 til 17.~september 2026, i alt 5.965 handelsdage. Perioden
begynder, hvor SHY har fuld historik.

Makrovariable fra FRED: rentekurvens hældning (\texttt{T10Y2Y}) og
arbejdsløsheden (\texttt{UNRATE}). Kurvehældningen er markedsdata og indgår med
samme dags værdi. Arbejdsløsheden offentliggøres først i måneden efter
referencemåneden, og observationsdatoen er derfor skubbet en måned frem, så
backtestet ikke handler på tal, der endnu ikke var kendt. Uden den justering
overvurderes strategiens præstation systematisk.

\section{Metode}
Regimeindikatoren er en sum af tre binære signaler: positiv kurvehældning,
faldende arbejdsløshed over 12 måneder, og indeks over sit eget 200-dages
gennemsnit. Ved score $\geq 2$ holdes den offensive position, ellers den
defensive. Positionen sættes på rebalanceringsdagen og anvendes fra dagen efter.

Strategien køres $k=1,\dots,21$ gange, hvor $k$ angiver hvilken handelsdag i
måneden der rebalanceres. Alle 21 varianter implementerer samme model og har
samme forventede afkast; forskellen mellem dem er udelukkende et arbitrært
implementeringsvalg. Maksimalt drawdown er beregnet på kumulerede logafkast.

\subsection{Forventede fortegn}
Fastlagt før estimation. Alle 21 varianter bør have samme forventede afkast.
Forventningen var, at den realiserede spredning alligevel ville være
sammenlignelig med strategiens merafkast over benchmark. Var spredningen lille,
ville konklusionen være, at strategien er robust --- også et resultat værd at
rapportere.

\section{Resultater}
Tabel~\ref{tab:udfald} viser fordelingen af udfald over de 21 varianter.

\begin{table}[htbp]\centering
\caption{Udfald på tværs af 21 rebalanceringsdage, 2003--2026}
\label{tab:udfald}\small
\begin{tabularx}{0.9\textwidth}{Xccccc}
\toprule
& Min & Median & Maks & Spredning & Passiv \\
\midrule
Sharpe            & 0,618 & 0,678 & 0,735 & 0,032 & 0,582 \\
Årligt afkast     & 0,088 & 0,096 & 0,105 & 0,004 & 0,108 \\
Maks.\ drawdown   & $-0{,}411$ & $-0{,}411$ & $-0{,}411$ & 0,000 & $-0{,}803$ \\
\bottomrule
\end{tabularx}
\vspace{0.4em}
\begin{minipage}{0.9\textwidth}\footnotesize
\textit{Noter:} Drawdown på kumulerede logafkast. Passiv er en position i SPY
gennem hele perioden.
\end{minipage}
\end{table}

Spændet i Sharpe er $0{,}735-0{,}618=0{,}117$. Medianvarianten slår den passive
benchmark med $0{,}678-0{,}582=0{,}096$. Forskellen mellem den heldigste og den
uheldigste implementering af den samme model er altså \emph{større} end hele
modellens fortrin over ikke at gøre noget. Afkastet er samtidig lavere end
passiv i alle 21 varianter, så Sharpe-gevinsten kommer udelukkende fra
risikoreduktion.

Det samlede drawdown er identisk for alle 21 varianter, hvilket først ligner en
fejl. Forklaringen er substantiel: alle varianter har samme peak
(19.~februar 2020) og samme trough (23.~marts 2020), fordi ingen af dem nåede
at gå defensivt under COVID-krakket. Regimescoren faldt fra 3 til 2, da indekset
brød sit 200-dages gennemsnit, men 2 er netop tærsklen, og de to
makrosignaler vendte ikke i tide.

Tabel~\ref{tab:episoder} opdeler derfor drawdown på de tre store nedture i
perioden.

\begin{table}[htbp]\centering
\caption{Maksimalt drawdown inden for hver episode}
\label{tab:episoder}\small
\begin{tabularx}{0.9\textwidth}{Xccccc}
\toprule
Episode & Min & Median & Maks & Spredning & Passiv \\
\midrule
Finanskrisen (okt.\ 2007--mar.\ 2009) & $-0{,}243$ & $-0{,}114$ & $-0{,}093$ & 0,046 & $-0{,}803$ \\
COVID (feb.--apr.\ 2020)              & $-0{,}411$ & $-0{,}411$ & $-0{,}411$ & 0,000 & $-0{,}411$ \\
Inflationsåret (jan.--okt.\ 2022)     & $-0{,}293$ & $-0{,}261$ & $-0{,}261$ & 0,007 & $-0{,}281$ \\
\bottomrule
\end{tabularx}
\end{table}

Tre ting følger. For det første leverer modellen kun beskyttelse i
finanskrisen, hvor tabet reduceres fra 80 til mellem 9 og 24 procent. For det
andet er netop den episode, hvor modellen virker, også den, hvor datovalget
betyder mest: forholdet mellem det værste og det bedste udfald er 2,6. For det
tredje gav modellen ingen beskyttelse i 2020 og var i 2022 i sit dårligste
udfald marginalt værre end passiv.

De 21 varianter holder samme position på 96,5 procent af dagene i gennemsnit,
og mindst 93,8 procent for det mest forskellige par. Antallet af regimeskift
over de 23 år er 15 til 19. Hele spredningen i udfald stammer altså fra et
tocifret antal beslutninger, hvor nogle varianter rammer få dage tidligere end
andre.

\section{Robusthed}
Tre tjek er gennemført. Episodeopdelingen ovenfor viser, at det samlede
drawdown skjuler et helt andet mønster end det, tabel~\ref{tab:udfald}
antyder. Positionsoverlappet dokumenterer, at varianterne er næsten identiske i
implementering og alligevel adskiller sig markant i udfald. Og optællingen af
regimeskift viser, at resultatet hviler på få begivenheder, hvilket i sig selv
begrænser, hvor stærkt spredningen kan tolkes.

Tre tjek udestår: \tbd{kvartalsvis rebalancering, hvor spredningen forventes
større}, \tbd{variation af regimetærsklen til $\geq 1$ og $\geq 3$} og
\tbd{en overlappende portefølje, der holder $1/21$ i hver variant, som et bud
på, hvad man kan gøre ved problemet}.

\section{Begrænsninger}
Regimemodellen er bevidst simpel og er ikke et forsøg på at replikere nogen
konkret forvalters model; resultatet siger noget om klassen af månedligt
rebalancerede regimestrategier, ikke om en bestemt af dem.
Offentliggørelsesforsinkelsen for arbejdsløsheden er approksimeret som præcis
én måned.

Vigtigst er dog, at perioden indeholder tre store nedture og 15--19
regimeskift. Spredningen mellem varianterne er dermed selv estimeret på få
begivenheder, og tallet $0{,}046$ for spredningen i finanskrisen hviler reelt
på én episode. Transaktionsomkostninger indgår ikke; med højst 19 skift over 23
år ville de dog ikke ændre konklusionen.

\section{Konklusion}
Et backtest af en månedligt rebalanceret regimestrategi rapporterer ét udfald
blandt 21 lige gyldige. På disse data spænder Sharpe over 0,117 på tværs af det
valg, mens strategiens fortrin over en passiv position er 0,096. Den
rapporterede præstation kan altså ikke adskilles fra held med en dato, ingen
har begrundet.

Konklusionen skærpes af episodeopdelingen. Modellens værdi ligger udelukkende i
finanskrisen, og netop dér svinger det undgåede tab med en faktor 2,6 mellem
den bedste og den værste rebalanceringsdag. I 2020, hvor nedturen var for
hurtig til, at makrosignalerne kunne nå at vende, gav modellen ingen
beskyttelse overhovedet. En strategi, der beskytter i én ud af tre kriser, og
hvis beskyttelse i det ene tilfælde varierer med en faktor 2,6 efter et
arbitrært implementeringsvalg, bør ikke rapporteres med ét Sharpe-tal.

\end{document}
